\chapter{Background and Related Work}
\label{chap:background}

\section{Transfer Learning and Full Fine-Tuning}

Transfer learning separates general representation learning from task adaptation. A pretrained
model learns linguistic regularities from a large corpus; supervised fine-tuning then adjusts it
to a labeled task. BERT demonstrated that a single bidirectional Transformer encoder could achieve
strong results across multiple language-understanding benchmarks after task-specific fine-tuning
\citep{devlin2019bert}. GLUE later provided a standardized collection of classification and
regression tasks for evaluating general language understanding \citep{wang2018glue}.

If a model has $P$ parameters, full fine-tuning stores gradients for approximately $P$ parameters
and, with an optimizer such as Adam, normally maintains additional first- and second-moment states.
The exact memory use depends on precision and implementation, but it scales with the full model
size. Repeating this process for $K$ tasks produces $K$ full checkpoints. This motivates methods
whose trainable and stored parameters are much smaller than $P$.

\section{Parameter-Efficient Fine-Tuning}

PEFT methods differ in where they introduce task-specific capacity.
Houlsby adapters insert small bottleneck networks into Transformer blocks
\citep{houlsby2019parameter}. Prefix tuning learns continuous vectors that act as virtual prefix
tokens for attention \citep{li2021prefix}. Prompt tuning learns a smaller set of input embeddings
\citep{lester2021power}. BitFit updates selected bias parameters \citep{zaken2022bitfit}.
These methods illustrate that effective task adaptation may occupy a lower-dimensional subspace
than the full parameter space.

PEFT is not automatically superior to full fine-tuning. Restricting trainable capacity introduces
an inductive bias, and the best method depends on model architecture, task, data volume, and
optimization. Its practical advantage lies in the reduced trainable state and the ability to keep
one shared base model.

\section{Low-Rank Adaptation}

Consider a pretrained matrix $W_0 \in \mathbb{R}^{d_{\mathrm{out}}\times d_{\mathrm{in}}}$.
LoRA keeps $W_0$ frozen and parameterizes its update as
\begin{equation}
  \Delta W = \frac{\alpha}{r}BA,
  \qquad
  A \in \mathbb{R}^{r\times d_{\mathrm{in}}},
  \quad
  B \in \mathbb{R}^{d_{\mathrm{out}}\times r}.
  \label{eq:lora-scale}
\end{equation}
The number of trainable parameters becomes
\begin{equation}
  P_{\mathrm{LoRA}} = r(d_{\mathrm{in}} + d_{\mathrm{out}})
  \label{eq:lora-params}
\end{equation}
instead of $d_{\mathrm{in}}d_{\mathrm{out}}$. When
$r \ll \defmin(d_{\mathrm{in}},d_{\mathrm{out}})$, the saving is substantial.
The scaling factor $\alpha/r$ controls the contribution of the learned update.

LoRA's main engineering advantage is predictability. A practitioner selects the rank, target
modules, dropout, and scaling factor, then trains using an ordinary optimization loop. However,
predictability comes from fixing the allocation. Every selected module usually receives the same
rank even when modules differ in size, function, or task relevance.

\subsection{The Rank-Selection Problem}

Rank is a capacity hyperparameter. Searching over ranks multiplies training cost. A rank selected
on one dataset or model may not transfer to another. DyLoRA trains a nested range of ranks so that
one run can support multiple deployment ranks \citep{valipour2023dylora}. QDyLoRA combines this
idea with quantized adaptation \citep{rajabzadeh2024qdylora}, while QLoRA focuses on reducing base
model memory through 4-bit quantization \citep{dettmers2023qlora}. These works reinforce that
rank and resource constraints are central practical questions rather than minor implementation
details.

\section{AdaLoRA}

AdaLoRA represents a weight update using a decomposition analogous to
\begin{equation}
  \Delta W = P\Lambda Q,
  \label{eq:adalora-decomp}
\end{equation}
where components associated with diagonal entries of $\Lambda$ can be scored and pruned.
It begins with a relatively generous rank and gradually reduces the total budget. Importance
reflects both parameter contribution and uncertainty, allowing capacity to move toward matrices
that appear more useful \citep{zhang2023adalora}.

The training schedule has three conceptual phases:
\begin{enumerate}
  \item an initial phase in which the larger parameterization is trained;
  \item a budget-reduction phase in which importance is periodically estimated and components are
  pruned;
  \item a final phase in which the resulting allocation is refined without further pruning.
\end{enumerate}

The PEFT implementation requires values corresponding to the initial rank, target rank, total
steps, initial warm-up, final fine-tuning period, and allocation interval. This makes AdaLoRA more
dependent on the training horizon than ordinary LoRA.

\section{Why AdaLoRA May See Less Practical Use}

This section explains technical barriers without claiming a comprehensive survey of adoption.

\subsection{Configuration Burden}

Fixed-rank LoRA has a compact mental model: choose $r$ and target modules. AdaLoRA adds multiple
interacting values. An unsuitable schedule can leave too few allocation updates in a short run or
can prune before the task signal is stable. The user must also ensure that the configured total
steps agree with the actual optimization loop.

\subsection{Noisy Importance in Short Runs}

Low-resource experiments may contain only tens or hundreds of optimizer steps. Early gradients
reflect initialization, minibatch order, and classifier-head instability. If an adaptive method
must allocate capacity during this period, its flexibility may amplify noise instead of finding a
better structure.

\subsection{Runtime and Implementation Complexity}

Dynamic allocation requires score maintenance and periodic mask or rank updates. The overhead may
be modest relative to very large models, but it matters in small experimental settings and makes
debugging harder. Fixed LoRA is supported by a broad ecosystem and has fewer moving parts.

\subsection{A Strong Simple Baseline}

Most importantly, fixed LoRA can already perform very well. A more complex allocator must beat not
only full fine-tuning but also a well-chosen fixed rank. The results in \Cref{chap:results} confirm
that this is a demanding baseline on SST-2.

\section{Dataset-Aware and Layer-Aware Adaptation}

Existing adaptive approaches mostly use signals observed during training, such as gradients,
weights, or singular values. Star-LoRA asks a complementary question: can useful allocation
decisions be made before training from properties of the downstream data? Pre-training analysis
cannot replace gradient evidence, but it can provide a prior. For example, a tiny imbalanced
dataset should not automatically receive the same capacity as a large balanced dataset.

The approach in this thesis is intentionally interpretable. It does not learn a rank controller
and does not claim that its handcrafted formula is globally optimal. Its purpose is to test whether
simple dataset evidence can improve a standard AdaLoRA configuration and to reveal where such a
policy fails.
